\documentclass[11pt]{amsart}

\usepackage[T1]{fontenc}
\usepackage{lmodern}
\usepackage{microtype}
\usepackage{amsmath,amssymb,amsthm}
\usepackage[hidelinks]{hyperref}

\hypersetup{
  pdftitle={A Counterexample to the Unrestricted Dargad--Larsson Domination Problem}
}

\newtheorem{theorem}{Theorem}
\newtheorem{corollary}[theorem]{Corollary}

\newcommand{\TS}{\mathcal{T}}

\title{A Counterexample to the Unrestricted Dargad--Larsson Domination Problem}
\date{}

\subjclass[2020]{91A46, 91A05}
\keywords{partizan subtraction game, truncated support, dominated option,
game order}

\begin{document}

\begin{abstract}
Dargad and Larsson asked in Problem~7.1 whether their Proposition~7.2
accounts for all dominated options in truncated-support games at the
knife's edge. We determine the complete order among the periodic game
values. Writing
\[
G_n=\TS_{b-a}(n;a,b),\qquad q=b-a+1,
\]
we prove, for $0\le r,s\le b$, that
\[
G_r\le G_s
\quad\Longleftrightarrow\quad
\left\lfloor\frac rq\right\rfloor=
\left\lfloor\frac sq\right\rfloor
\ \text{and}\ r\le s.
\]
Thus the values form increasing chains on consecutive $q$-blocks and
are incomparable across blocks. This confirms the contextual reading of
Problem~7.1 for the $b+1$ representative positions. Under its
unrestricted wording, the comparisons in the source are complete if and
only if $q\mid(a-1)$. Otherwise the statement is false: for
$(a,b)=(3,5)$, the Left option $G_4$ of $G_6$ is dominated by $G_5$, a
comparison outside the range of Proposition~7.2.
\end{abstract}

\maketitle

\section{Introduction}

In a partizan subtraction game, Left and Right remove tokens from a
single heap according to different subtraction sets; see
\cite{DHNP}. Dargad and Larsson introduced the truncated-support game
$\TS_\tau(n;a,b)$, in which
\[
S_L=\{1,\dots,a\},\qquad S_R=\{\tau+1,\dots,b\},
\qquad 0<a<b.
\]
At the knife-edge truncation $\kappa=b-a$, they proved the game-value
periodicity
\begin{equation}\label{eq:periodicity}
\TS_\kappa(n;a,b)=\TS_\kappa(t;a,b),
\qquad t\equiv n\pmod{b+1},\quad 0\le t\le b,
\end{equation}
and asked whether the only dominated options are those given by their
Proposition~7.2 \cite[Problem~7.1]{DL}; that proposition is stated
there under the hypothesis $3\le a<b$.

The paragraph preceding that proposition restricts attention to the
representatives with heap sizes $0,\dots,b$, whose options have heap
sizes at most $b-1$. This is why the comparison range of
Proposition~7.2, quoted at the end of this section, ends at $b-2$: the
last consecutive pair of option heap
sizes available to a representative is $(b-2,b-1)$. Heap size $b+1$ is
the first position with an option of heap size $b$, and the residue $b$
is the one the source's comparisons never reach. The printed
question itself contains no restriction on $n$. We resolve both
readings. The structural result is the exact
partial order on the $b+1$ periodic values; the unrestricted reading
then has the counterexample stated in the abstract.

Put
\[
q=b-a+1,\qquad G_n=\TS_{b-a}(n;a,b).
\]
At the knife edge, Left removes an element of $\{1,\dots,a\}$ and Right
removes an element of $\{q,\dots,b\}$. Hence, for $0\le n\le b$,
\begin{equation}\label{eq:recursion}
G_n=
\left\{
G_{n-i}:1\le i\le\min(a,n)
\ \middle|\
G_{n-j}:q\le j\le\min(b,n)
\right\}.
\end{equation}

\subsection*{Statements from the source}

We reproduce the three statements of \cite{DL} on which the discussion
below turns; the only changes are the symbol $\TS$ for the ruleset and
ordinary set notation. Observation~5.2 reads: ``Let $a,b$ and $\tau$ be
positive integers such that $a<b$ and $\tau<b$. If $n\le\tau$, then
Right does not have any move, whereas Left can remove tokens one by
one. Thus, for all $1\le n\le\tau$, $\TS_\tau(n;a,b)=n$.''
Proposition~7.2 reads: ``Consider $G=\TS_\kappa(n;a,b)$ and
$H=\TS_\kappa(n+1;a,b)$ where $3\le a<b$. Then, $G\le H$ for all
$n\in\{b-a+1,\dots,b-3,b-2\}\setminus Q$'', where
\[
Q=\{(m+1)(b-a)+m:m\in\mathbb{Z}_{\ge1}\}=\{(m+1)q-1:m\ge1\}
\]
is the set introduced in the paragraph preceding it. That paragraph
presents the games $\TS_\kappa(n;a,b)$ with $0\le n\le b$ as
representatives of the canonical forms, and records: ``The options of
these games have heap sizes ranging from $0$ to $b-1$. Thus, only the
games with heap sizes $0,1,\dots,b-1$ are required to be compared.''
Problem~7.1 reads:
``Is it true that the only dominated options of $\TS_{b-a}(n;a,b)$ are
given by Proposition~7.2.''

\section{The order of the periodic values}

\begin{theorem}\label{thm:order}
For all $0\le r,s\le b$,
\begin{equation}\label{eq:order}
G_r\le G_s
\quad\Longleftrightarrow\quad
\left\lfloor\frac rq\right\rfloor=
\left\lfloor\frac sq\right\rfloor
\ \text{and}\ r\le s.
\end{equation}
Consequently, if
\[
B_m=\{mq,mq+1,\dots,\min((m+1)q-1,b)\},
\]
then the order on $G_0,\dots,G_b$ is the disjoint union of the strict
chains indexed by the nonempty blocks $B_m$. In particular, these
$b+1$ values are pairwise distinct, and $b+1$ is the least positive
period in \eqref{eq:periodicity}.
\end{theorem}

\begin{proof}
We use strong induction on $r+s$. The case $r=s=0$ is immediate. Recall
that, for short games, $X\le Y$ if and only if no Left option $X^L$
satisfies $Y\le X^L$ and no Right option $Y^R$ satisfies $Y^R\le X$
\cite{Siegel}.

Suppose first that $r$ and $s$ lie in the same $q$-block and $r\le s$.
Every Left option $G_u$ of $G_r$ has $u<r\le s$, so the induction
hypothesis gives $G_s\nleq G_u$. Every Right option $G_v$ of $G_s$
satisfies $v\le s-q$ and therefore lies in a block strictly before the
common block of $r$ and $s$; in the first block there are no Right
options. Hence $G_v\nleq G_r$ by induction, and the comparison criterion
gives $G_r\le G_s$.

If $r$ and $s$ lie in the same block and $r>s$, then $G_{r-1}$ is a
Left option of $G_r$, and $r-1$ remains in that block. By induction,
$G_s\le G_{r-1}$, so $G_r\nleq G_s$.

Now suppose that the block of $r$ is later than the block of $s$, and
set
\[
u=\max\{s,r-a\}.
\]
Since $r\le b=a+q-1$, we have $r-a\le q-1$. If $s<q$, then $u<q$;
if $s\ge q$, then $u=s$. Thus $u$ lies in the same block as $s$.
Moreover, $r-a\le u<r$, so $G_u$ is a Left option of $G_r$. Induction
gives $G_s\le G_u$, and therefore $G_r\nleq G_s$.

Finally, suppose that the block of $r$ is earlier than the block of
$s$, and set
\[
v=\min\{r,s-q\}.
\]
Writing $i=\lfloor r/q\rfloor$ and $j=\lfloor s/q\rfloor$, we have
$j>i$ and
\[
s-q\ge(j-1)q\ge iq,
\qquad r<(i+1)q.
\]
Hence $v$ lies in the same block as $r$. Also $q\le s-v\le b$, so
$G_v$ is a Right option of $G_s$. By induction, $G_v\le G_r$, and
therefore $G_r\nleq G_s$. This proves \eqref{eq:order}.
\end{proof}

\section{Domination and the counterexample}

For a heap size $n$, reduce each option's target heap size modulo
$p=b+1$ to a residue in $\{0,\dots,b\}$. Distinct options of the same
player have distinct residues: the difference between two legal move
sizes is at most $a-1<p$. Theorem~\ref{thm:order} therefore gives the
complete domination rule.

\begin{corollary}\label{cor:domination}
Among Left options whose target residues lie in one $q$-block, every
option except the one with largest residue is dominated. Among Right
options whose target residues lie in one $q$-block, every option except
the one with smallest residue is dominated. No domination occurs
between options whose target residues lie in different blocks.
\end{corollary}

Throughout, ``given by Proposition~7.2'' means ``derivable from
Observation~5.2 and Proposition~7.2 by transitivity and the periodicity
\eqref{eq:periodicity}''. This is the reading intended by the paragraph
preceding Proposition~7.2, and both verdicts below are
taken under it.
Without that latitude the answer is negative for every $a\ge2$, and for
a trivial reason: the comparisons inside the first block come from
Observation~5.2 and not from Proposition~7.2, whose range begins at $q$.

Recall that $Q=\{(m+1)q-1:m\ge1\}$ is the exceptional set of
Proposition~7.2. For $q\le n\le b-1$, Theorem~\ref{thm:order} gives
\[
G_n<G_{n+1}\quad\Longleftrightarrow\quad n\notin Q;
\]
when $n\in Q$, the two values are incomparable. Thus, for $3\le a<b$,
Proposition~7.2 of \cite{DL} supplies all consecutive comparisons
required among the option values $G_0,\dots,G_{b-1}$ of the
representative positions, while Observation~5.2 supplies the first
block. This answers the
contextual, representative-position reading of Problem~7.1
affirmatively.

For the unrestricted reading, the omitted residue $b$ is decisive.
Since
\[
b=q+a-1,
\]
the final block is the singleton $\{b\}$ exactly when $q\mid(a-1)$.
In that case Corollary~\ref{cor:domination} gives no comparison involving
$G_b$, so Observation~5.2 and Proposition~7.2, together with
transitivity and periodicity, account for every dominated option at
every heap size.

If $q\nmid(a-1)$, then $b-1$ and $b$ lie in the same block, and
Theorem~\ref{thm:order} gives
\[
G_{b-1}<G_b.
\]
Since $q\mid0$, the hypothesis $q\nmid(a-1)$ forces $a\ge2$, so both
values are Left options of $G_{b+1}$, reached by removing $2$ and $1$
tokens, respectively. Hence the option $G_{b-1}$ is dominated by
$G_b$. The comparison occurs at index $b-1$, outside the printed range
$n\le b-2$ of Proposition~7.2. Thus, under the convention above, the
unrestricted statement holds if and only if $q\mid(a-1)$.

The two values of $a$ outside the hypothesis of Proposition~7.2 agree
with this. For $a=1$ each player has a single legal move size, so no
position has two options of the same player and nothing is dominated;
the statement is true, and indeed $q=b$ divides $a-1=0$. For $a=2$ the
range $\{q,\dots,b-2\}$ of Proposition~7.2 is empty, so the proposition
supplies nothing, while Observation~5.2 covers exactly the first block
$B_0=\{0,\dots,b-2\}$; since $q=b-1$ does not divide $a-1=1$, the
comparison $G_{b-1}<G_b$ is again unsupplied and the statement is
false.

In particular, for $(a,b)=(3,5)$ one has $q=3$ and
\[
G_3<G_4<G_5.
\]
At heap size $6$, Left may move to $G_5,G_4$, or $G_3$, so $G_4$ is
dominated by $G_5$. Proposition~7.2 supplies only the comparison
$G_3\le G_4$ for this ruleset. Thus the unrestricted form of
Dargad and Larsson's Problem~7.1 is false.

\begin{thebibliography}{9}

\bibitem{DHNP}
E.~Duch\^ene, M.~Heinrich, R.~J.~Nowakowski, and A.~Parreau,
\emph{Partizan subtraction games},
Integers \textbf{21B} (2021), Paper A8.

\bibitem{DL}
A.~Dargad and U.~Larsson,
\emph{Partizan subtraction with full and truncated support},
arXiv:2607.27989v1 [math.CO], 2026.

\bibitem{Siegel}
A.~N.~Siegel,
\emph{Combinatorial Game Theory},
Graduate Studies in Mathematics, vol.~146,
American Mathematical Society, Providence, RI, 2013.

\end{thebibliography}

\end{document}